\documentclass{situationist-preprint}

\title{Incongruity and Its Resolution}
\msid{AES-13-a}
\msfamily{AES-13}
\msdiscipline{Aesthetics}
\mstype{Principal manuscript}
\mskeywords{humor; jokes; incongruity theory; resolution; amusement; superiority theory}

\begin{document}
\maketitle

\begin{abstract}
What makes a joke funny? Incongruity theories hold that amusement is a response to the perception of something incongruous, but bare incongruity is often merely puzzling, and incongruity-resolution theories hold that amusement arises when an incongruity is resolved, but a fully resolved incongruity is a solved puzzle. We argue that a joke is funny when it presents an incongruity that is resolved, but not entirely: the punch line supplies a reading on which the incongruous element fits, and that reading is itself in tension with what we know. Amusement is the response to holding both at once. We show that the account explains why unresolved incongruity is not funny, why fully resolved incongruity is not funny, and why jokes cease to be funny when they are explained; and we argue that it subsumes what is right in superiority and release theories.
\end{abstract}

\section{The problem}

\citet{Hobbes1651} held that laughter expresses a sudden glory arising from the perception of some eminency in ourselves by comparison with the infirmity of others (\emph{Leviathan} I.6). \citet{Hutcheson1750} objected that we laugh at much that involves no comparison with ourselves, and located the source of laughter in the contrast of ideas. Kant described laughter as an affect arising from the sudden transformation of a strained expectation into nothing \citep[§54]{Kant1790}, and Schopenhauer as the sudden perception of an incongruity between a concept and the real objects thought through it (\emph{World as Will and Representation} I, §13). \citet{Freud1905} explained the pleasure of jokes as a saving of psychic expenditure, and in particular of the energy otherwise spent on inhibition.

The incongruity tradition has dominated recent work \citep{Morreall1983,Carroll2014}, but it faces a problem stated by the psychologist \citet{Suls1972}: incongruity alone is not sufficient for humor, since many incongruities are puzzling, alarming or merely nonsensical. Suls proposed a two-stage model in which the listener first meets an incongruity and then finds a cognitive rule that resolves it. \citet{Oring1992} argued that the resolution typically reveals an \emph{appropriate} incongruity, a relation that is appropriate in one respect and inappropriate in another. \citet{HurleyDennettAdams2011} proposed that humor rewards the detection of mistaken commitments in the mental spaces we construct while understanding, and \citet{McGrawWarren2010} that it arises from violations perceived as benign.

\section{The account}

\begin{definition}[Incongruity]\label{d:inc}
An element of a joke is \emph{incongruous} if it conflicts with the expectations that the joke's setup has led its audience to form.
\end{definition}

\begin{definition}[Partial resolution]\label{d:part}
An incongruity is \emph{resolved} if the audience finds a reading of the joke on which the incongruous element fits. It is resolved \emph{partially} if that reading is itself in tension with what the audience knows or believes about the world, so that it makes sense of the element without making it unremarkable.
\end{definition}

\begin{principle}[Amusement]\label{p:amuse}
Amusement at a joke is the pleasure of grasping a resolution of an incongruity while continuing to perceive what remains incongruous in the resolution.
\end{principle}

Consider Groucho Marx's line: \enquote{One morning I shot an elephant in my pajamas. How he got into my pajamas I'll never know.} The first sentence establishes the reading on which the speaker wore the pajamas. The second is incongruous with it. The incongruity is resolved: there is a parse of the first sentence, available all along, on which the elephant was in the pajamas. But the resolution is in tension with what we know of elephants and pajamas, and the speaker's puzzlement treats the absurd reading as a fact to be explained. The joke is funny because the reanalysis both succeeds and fails.

\begin{theorem}\label{t:main}
A joke is funny when, and because, it presents an incongruity that it resolves partially: a joke whose incongruity is not resolved, or is resolved entirely, is not funny.
\end{theorem}

\begin{proof}
By Principle~\ref{p:amuse}, amusement requires both a resolution and a residual incongruity. If the incongruity is not resolved, the audience has no reading on which the element fits, and the response is puzzlement: the joke is nonsense. If it is resolved entirely, the resolving reading is in no tension with what the audience knows, and the element becomes unremarkable: the audience has solved a riddle or learned a fact, and the response is the satisfaction of understanding. Only when the resolution is partial (Definition~\ref{d:part}) are both conditions met, and the joke amuses.
\end{proof}

A corollary concerns the explained joke, discussed in the companion manuscript AES-13-b. An explanation makes the resolving reading explicit and accounts for its tension with what we know by naming it as a device. The residual incongruity is thereby itself resolved, and the joke ceases to amuse: what was held in tension has been put in order.

\section{Objections and replies}

\begin{objection}[Superiority]\label{o:sup}
Many jokes amuse by making someone look foolish, which the incongruity account does not capture \citep{Hobbes1651}.
\end{objection}

\begin{reply}
The foolishness in such jokes is a partially resolved incongruity: the butt of the joke acts on a reading of the situation that makes sense from his point of view and is absurd from ours. The pleasure of superiority, where it is present, accompanies amusement but is not necessary for it, as wordplay and absurdist jokes, which disparage no one, show.
\end{reply}

\begin{objection}[Benign violation]\label{o:bv}
What matters is not resolution but that the violation be perceived as harmless \citep{McGrawWarren2010}.
\end{objection}

\begin{reply}
Benignity is a condition on amusement and not a rival account of it. A partially resolved incongruity whose content is too distressing will not amuse, because distress crowds out the pleasure Principle~\ref{p:amuse} describes. But benign violations are not all funny: a harmless breach of etiquette, fully understood, is merely a breach. Resolution and its incompleteness explain what benignity leaves open.
\end{reply}

\section{Conclusion}

A joke is funny when it resolves an incongruity, but not entirely. The resolution distinguishes a joke from nonsense; its incompleteness distinguishes it from a solved puzzle. Amusement is the pleasure of holding the resolution and its residue together, which is why it does not survive an explanation that dissolves the residue.

\begin{thebibliography}{99}\small
\bibitem[Carroll(2014)]{Carroll2014} Carroll, N. (2014). \emph{Humour: A Very Short Introduction}. Oxford University Press.
\bibitem[Freud(1905)]{Freud1905} Freud, S. (1905). \emph{Der Witz und seine Beziehung zum Unbewussten}. Deuticke. Trans. J. Strachey as \emph{Jokes and Their Relation to the Unconscious}.
\bibitem[Hobbes(1651)]{Hobbes1651} Hobbes, T. (1651). \emph{Leviathan}. Andrew Crooke.
\bibitem[Hurley et~al.(2011)]{HurleyDennettAdams2011} Hurley, M.~M., Dennett, D.~C. and Adams, R.~B. (2011). \emph{Inside Jokes: Using Humor to Reverse-Engineer the Mind}. MIT Press.
\bibitem[Hutcheson(1750)]{Hutcheson1750} Hutcheson, F. (1750). \emph{Reflections upon Laughter, and Remarks upon the Fable of the Bees}. R. Urie.
\bibitem[Kant(1790)]{Kant1790} Kant, I. (1790). \emph{Critique of the Power of Judgment}. Trans. P. Guyer and E. Matthews. Cambridge University Press, 2000.
\bibitem[McGraw and Warren(2010)]{McGrawWarren2010} McGraw, A.~P. and Warren, C. (2010). Benign violations: Making immoral behavior funny. \emph{Psychological Science}, 21(8), 1141--1149.
\bibitem[Morreall(1983)]{Morreall1983} Morreall, J. (1983). \emph{Taking Laughter Seriously}. State University of New York Press.
\bibitem[Oring(1992)]{Oring1992} Oring, E. (1992). \emph{Jokes and Their Relations}. University Press of Kentucky.
\bibitem[Suls(1972)]{Suls1972} Suls, J.~M. (1972). A two-stage model for the appreciation of jokes and cartoons: An information-processing analysis. In J.~H. Goldstein and P.~E. McGhee (Eds.), \emph{The Psychology of Humor: Theoretical Perspectives and Empirical Issues}. Academic Press.
\end{thebibliography}

\end{document}
